\chapter{From one polylogarithm to many}
\label{ch:01-foundations}

An ordinary polylogarithm has two descriptions: a series and an iterated
integral. Nesting the first and extending the alphabet of the second produce
the same larger class of functions. Their two multiplication laws then provide
exact relations. This chapter fixes conventions before those relations are
used to organize special values.

\section{Multiple polylogarithms: two faces}

The classical polylogarithm $\Li_s(z)=\sum_{n\ge1}z^n/n^s$ has, for the reader of Lewin, two complementary descriptions: a power series and an iterated integral $\Li_s(z)=\int_0^z\Li_{s-1}(t)\,dt/t$. Each survives the passage to many indices, and the resulting two descriptions of one and the same function are the source of essentially every nontrivial identity in this circle of ideas. We set up both here; the products they generate occupy the next section.

\subsection{The nested sum}

\begin{definition}
For $s_1,\dots,s_k\in\mathbb{C}$ and $z_1,\dots,z_k\in\mathbb{C}$ the \emph{multiple polylogarithm} is the strictly nested sum
\[
  \Li_{s_1,\dots,s_k}(z_1,\dots,z_k)
  \;=\;\sum_{n_1>n_2>\dots>n_k\ge1}\;\prod_{i=1}^{k}\frac{z_i^{\,n_i}}{n_i^{\,s_i}} .
\]
Its \emph{weight} is $w=s_1+\dots+s_k$ and its \emph{depth} is $k$. The single-index case $k=1$ recovers $\Li_s(z)$.
\end{definition}

Two features of this definition demand comment before anything is proved.

\begin{remark}[Conventions differ; ours is fixed]\label{found:rem:conv}
The inequality is \emph{strict}, $n_1>n_2>\dots>n_k$. A sizable part of the literature (and several computer-algebra systems) instead uses a non-strict summation $n_1\ge n_2\ge\dots\ge n_k$, or reverses the role of the indices so that the \emph{innermost} summation variable carries the first parameter. These choices are genuinely different functions (they differ by lower-depth terms, exactly the diagonal contributions $n_i=n_{i+1}$), and the discrepancy is the single most common source of sign and indexing confusion when comparing sources. We adhere throughout to the strict, ``outermost variable carries $s_1$'' convention above. The arguments are likewise ordered with $z_1$ paired to the largest index $n_1$. When checking against \texttt{MultiplePolyLog}, confirm the convention numerically on a small case rather than assuming agreement.
\end{remark}

\begin{remark}[Multiple zeta values]
Setting every $z_i=1$ gives the \emph{multiple zeta value}
\[
  \zeta(s_1,\dots,s_k)=\Li_{s_1,\dots,s_k}(1,\dots,1)
  =\sum_{n_1>\dots>n_k\ge1}\prod_i n_i^{-s_i},
\]
the central special values of the theory (Zagier 1994; Hoffman 1992). The strict nesting is what makes $\zeta(s_1,\dots,s_k)$ finite for $s_1\ge2$.
\end{remark}

\paragraph{Convergence.} The single-variable series $\Li_s(z)$ converges absolutely for $|z|<1$ (any $s$), on $|z|=1$ for $\Re s>1$, and conditionally on $|z|=1$, $z\ne1$, for $\Re s>0$. For the nested sum the corresponding statement is clean.

\begin{proposition}[A sufficient absolute-convergence domain]
For positive integer indices, the nested series converges absolutely when
$|z_i|\le1$ for all $i$ and $s_1\ge2$. It also converges absolutely for
$|z_1|<1$, $|z_i|\le1$ for $i>1$, with $s_1\ge1$.
\end{proposition}
\begin{proof}
The absolute value of the inner nested sum is at most
$H_{n_1-1}^{k-1}=O((1+\log n_1)^{k-1})$.
The outer series is therefore bounded by
$\sum |z_1|^{n_1}(1+\log n_1)^{k-1}/n_1^{s_1}$,
which converges in either stated case.
\end{proof}
These are sufficient conditions, not a complete classification of every
boundary point. Oscillation can give conditional convergence: $\Li_1(-1)$
converges, but $\sum |(-1)^n/n|$ diverges. More generally, away from the
polydisc an interior domain is given by
$|z_1\cdots z_j|<1$ for each $1\le j\le k$.
At unity, for positive integer indices, admissibility is exactly $s_1\ge2$.

\subsection{The iterated integral}

The second face is Chen's iterated integral. Fix the two differential one-forms
\[
  \omega_0=\frac{dt}{t},\qquad \omega_1=\frac{dt}{1-t},
\]
and, for a word $u=a_1a_2\cdots a_w$ in the letters $\{0,1\}$, write the iterated integral over the standard simplex
\[
  \int_0^z u \;=\; \int_{0<t_1<t_2<\dots<t_w<z}\ \omega_{a_w}(t_1)\,\omega_{a_{w-1}}(t_2)\cdots\omega_{a_1}(t_w),
\]
i.e.\ the leftmost letter $a_1$ is attached to the \emph{outermost} (largest) variable $t_w$ and one integrates from the top down. This is the orientation under which $d/dz$ strips the outermost form. With this normalization the one-variable multiple polylogarithm is an iterated integral over a single explicit word.

\begin{theorem}[Iterated-integral representation]\label{found:thm:iter}
For $s_1\ge2$ and the remaining $s_i\ge1$,
\[
  \Li_{s_1,\dots,s_k}(z,1,\dots,1)
  \;=\;\int_0^z\ \underbrace{0^{\,s_1-1}1}_{\text{block }1}\,\underbrace{0^{\,s_2-1}1}_{\text{block }2}\cdots\underbrace{0^{\,s_k-1}1}_{\text{block }k},
\]
the iterated integral of the binary word
\[
  u(\mathbf s)=0^{\,s_1-1}1\,0^{\,s_2-1}1\,\cdots\,0^{\,s_k-1}1
\]
of length $w$. Equivalently, with the alternative kernel $\widetilde\omega_1=dt/(t-1)=-\omega_1$ one carries one sign per $1$-letter, giving the textbook normalization
\[
  \Li_{s_1,\dots,s_k}(z,1,\dots,1)=(-1)^k\!\!\int_{0<t_1<\dots<t_w<z}\!\!\prod_{j=1}^{w}\eta_{a_{w+1-j}}(t_j),\qquad
  \eta_0=\tfrac{dt}{t},\ \eta_1=\tfrac{dt}{t-1},
\]
the factor $(-1)^k$ being exactly the number of $1$-letters in $u(\mathbf s)$, one per block.
\end{theorem}

\paragraph{The composition $\leftrightarrow$ word dictionary.} The map $\mathbf s=(s_1,\dots,s_k)\mapsto u(\mathbf s)$ is a bijection between compositions (ordered tuples with each $s_i\ge1$) and binary words that \emph{end in a $1$}. Reading $u$ from left to right: a $1$ \emph{closes a block} and contributes a new summation index; the run of $0$'s immediately preceding that $1$ sets the corresponding exponent, $s_i-1$ zeros for exponent $s_i$. Thus depth $k=$ (number of $1$'s) and weight $w=$ (length of the word). Admissibility $s_1\ge2$ is the requirement that $u$ \emph{begin with a $0$}: a leading $1$ would attach $\omega_1$ to the top variable and produce the logarithmic divergence at $t_w\to1$ when $z=1$. The endpoints play dual roles --- $\omega_0$ is singular at $0$, $\omega_1$ at $1$ --- so a convergent word neither starts with $1$ nor ends with $0$.

\begin{proof}[Origin]
Expand the innermost form geometrically. The depth-one case is Lewin's
$\Li_s(z)=\int_0^z 0^{s-1}1$, obtained by iterating $\Li_s(z)=\int_0^z\Li_{s-1}(t)\,\omega_0$ down to $\Li_1(z)=-\log(1-z)=\int_0^z\omega_1$. For general depth, $\int_0^z\omega_1$ on the innermost block introduces $\sum_{n_k\ge1}t^{n_k}/n_k$, the subsequent $\omega_0$'s raise the power of $n_k$, and the next $\omega_1$ couples a strictly larger index $n_{k-1}>n_k$; the strict inequality is forced by the nesting of the simplex $t_1<\dots<t_w$. Carrying this out reproduces the nested sum, and Chen's theory (1977) guarantees convergence and homotopy invariance of the integral. The general framework is due to Goncharov; the systematic exploitation for these functions is Goncharov's and Borwein--Bradley--Broadhurst--Lisone\v{k}'s.
\end{proof}

\subsection{General alphabets and a fixed integral orientation}
For a path from $0$ to $z$ avoiding its letters, use the outer-letter-first
recursion
\[
G(a_1,\ldots,a_w;z)=\int_0^z\frac{G(a_2,\ldots,a_w;t)}{t-a_1}\,dt,
\qquad G(;z)=1.
\]
The all-zero word is regularized as $G(0^w;z)=\log^w z/w!$.
For positive integer indices the dictionary in our decreasing-index
convention is
\[
\Li_{s_1,\ldots,s_k}(z_1,\ldots,z_k)
=(-1)^kG\left(0^{s_1-1},z_1^{-1},
 0^{s_2-1},(z_1z_2)^{-1},\ldots,
 0^{s_k-1},(z_1\cdots z_k)^{-1};1\right).
\]
Thus the letters are inverse \emph{prefix} products. Reversing the
integration order also reverses the letter word; mixing these conventions
changes the function. The formula follows by expanding each nonzero kernel
geometrically near the origin and integrating recursively, then continuing
along a common path. Wolfram's generalized-polylogarithm interface must be
calibrated against this orientation before numerical comparisons.

\subsection{Two facts to carry forward}

\begin{enumerate}
\item[(1)] \textbf{The differential strips the first letter.} In the integral orientation above,
\[
  z\frac{d}{dz}\,\Li_{s_1,\dots}(z,1,\dots)=\begin{cases}\Li_{s_1-1,s_2,\dots}(z,1,\dots),& s_1\ge2,\\[2pt]
  \dfrac{z}{1-z}\,\text{(remove the leading block)},& s_1=1,\end{cases}
\]
i.e.\ differentiation removes one outer letter --- a $0$ lowers $s_1$, the leading $1$ deletes the first index. This recursion (equivalently the integral equation defining the word) drives every analytic continuation and every series-rearrangement proof. It is also the recursion behind the appearance of \texttt{HarmonicPolyLog} in \texttt{DSolve}/\texttt{AsymptoticDSolveValue} output for Fuchsian systems with rational coefficients (Remiddi--Vermaseren 1999 for the $\{0,1\}$-alphabet case).
\item[(2)] \textbf{One function, two presentations.} The \emph{same} $\Li_{s_1,\dots,s_k}$ is simultaneously the nested sum and the iterated integral. Multiplying two nested sums and re-sorting the indices yields the \emph{stuffle} (harmonic, quasi-shuffle) product $\ast$; multiplying two iterated integrals over a common path and shuffling the letters yields the \emph{shuffle} product $\shuffle$. Both express a product of polylogarithms as a $\mathbb{Q}$-linear combination of polylogarithms, but through different combinatorics; equating them is the engine of the next section.
\end{enumerate}

\subsection{The smallest nontrivial example, in full}

Take depth $k=2$, $\mathbf s=(2,1)$, the lightest admissible non-classical case. The nested sum is
\[
  \Li_{2,1}(z,1)=\sum_{m>n\ge1}\frac{z^m}{m^2\,n}
  =\sum_{m=2}^{\infty}\frac{z^m}{m^2}\,H_{m-1},\qquad H_{m-1}=\sum_{n=1}^{m-1}\tfrac1n,
\]
exhibiting the characteristic harmonic-number coefficient. The associated word is
\[
  u(2,1)=\underbrace{0^{\,2-1}1}_{\text{block }(s_1=2)}\ \underbrace{0^{\,1-1}1}_{\text{block }(s_2=1)}=011,
\]
a word of length $w=3$ (weight $3$) with two $1$'s (depth $2$), beginning in $0$ (admissible) and ending in $1$ (convergent). By Theorem~\ref{found:thm:iter},
\[
  \Li_{2,1}(z,1)=\int_0^z 011
  =\int_{0<t_1<t_2<t_3<z}\frac{dt_3}{t_3}\,\frac{dt_2}{1-t_2}\,\frac{dt_1}{1-t_1},
\]
the leading $0$ sitting on the outermost variable $t_3$ and the two $1$'s on $t_2,t_1$. (Equivalently $(-1)^2=+1$ in the $\eta_1=dt/(t-1)$ normalization, consistent with the unsigned $\omega_1=dt/(1-t)$ form.) A numerical check at $z=0.7$ gives both sides equal to $0.2384965\ldots$, and at $z=1$ the integral $\int_0^1 011$ recovers Euler's celebrated
\[
  \zeta(2,1)=\sum_{m>n\ge1}\frac{1}{m^2 n}=\zeta(3),
\]
the depth-two value whose collapse to a depth-one zeta was Euler's first multiple-zeta identity --- and, in Wolfram~15, exactly the symbolic output of \texttt{MultipleZeta} on $\{2,1\}$.


\section{Multiple zeta values and their relations}

Specializing the multiple polylogarithm at $z_1=\dots=z_k=1$ produces the central objects of this chapter, the \emph{multiple zeta values} (MZVs):
\begin{equation}
  \zeta(s_1,\dots,s_k) \;=\; \Li_{s_1,\dots,s_k}(1,\dots,1)
  \;=\; \sum_{n_1>n_2>\dots>n_k\ge1}\;\prod_{i=1}^{k}\frac{1}{n_i^{\,s_i}}.
  \label{found:eq:mzv-def}
\end{equation}
The outermost summation index $n_1$ carries the exponent $s_1$, and the strict inequalities make the sum a genuine iterated, rather than ordinary, multiple series. The series \eqref{found:eq:mzv-def} converges if and only if $s_1\ge2$; such an index tuple is called \emph{admissible}. (Any $s_i\ge1$ for $i\ge2$ is allowed; only the leading entry is constrained.) We call $w=s_1+\dots+s_k$ the \emph{weight} and $k$ the \emph{depth}. In Wolfram Language these are \texttt{MultipleZeta}, with $\zeta(2,1)$ written \texttt{MultipleZeta[\{2,1\}]}.

The depth-one values are the classical zeta values. Euler's formula is
\[
\zeta(2n)=\frac{(-1)^{n+1}B_{2n}(2\pi)^{2n}}{2(2n)!}
\in\pi^{2n}\mathbb Q.
\] The first genuinely new phenomenon is that distinct MZVs can collapse onto a single lower-depth value. The prototype is Euler's identity
\begin{equation}
  \zeta(2,1)\;=\;\sum_{n>m\ge1}\frac{1}{n^2\,m}\;=\;\zeta(3),
  \label{found:eq:euler-21}
\end{equation}
already known to Euler in the 1740s. The conjecture is that the $\mathbb{Q}$-vector space spanned by all MZVs is \emph{weight-graded}: products and relations never mix weights, so one studies each weight $w$ separately.

\subsection*{The sum theorem}

Euler's $\zeta(2,1)=\zeta(3)$ is the depth-two case of the first structural theorem. Summing all admissible MZVs of a fixed weight $w$ and fixed depth $k$ recovers the single zeta value $\zeta(w)$.
\begin{theorem}[Sum theorem; Euler for $k=2$, Granville and Zagier in general, 1990s]
For all $w>k\ge1$,
\begin{equation}
  \sum_{\substack{s_1+\dots+s_k=w\\ s_1\ge2,\;s_i\ge1}}\zeta(s_1,\dots,s_k)\;=\;\zeta(w).
  \label{found:eq:sum-thm}
\end{equation}
\end{theorem}
\noindent At $w=3,k=2$ the only admissible tuple is $(2,1)$, recovering \eqref{found:eq:euler-21}. At $w=4,k=2$ the admissible tuples are $(3,1)$ and $(2,2)$, so
\begin{equation}
  \zeta(3,1)+\zeta(2,2)=\zeta(4)=\tfrac{\pi^4}{90}.
  \label{found:eq:sum-wt4}
\end{equation}
This one equation does not separate $\zeta(3,1)$ from $\zeta(2,2)$; for that we need a \emph{second} multiplication law, to which we now turn.

\subsection*{The two products}

MZVs admit two structurally different multiplication rules, because $\zeta(s_1,\dots,s_k)$ has two faithful representations: as a nested sum \eqref{found:eq:mzv-def}, and as an iterated integral. Each representation makes the product of two MZVs expand into a sum of MZVs, but the two expansions differ --- and that discrepancy is the engine of the theory.

\paragraph{(a) The stuffle (harmonic / quasi-shuffle) product.}
Multiply two nested sums directly. For depth one,
\[
  \zeta(a)\,\zeta(b)=\Bigl(\sum_{m\ge1}m^{-a}\Bigr)\Bigl(\sum_{n\ge1}n^{-b}\Bigr)
  =\sum_{m>n}+\sum_{n>m}+\sum_{m=n},
\]
splitting the lattice $\{(m,n)\}$ by the trichotomy $m>n$, $m<n$, $m=n$. The three regions are exactly two depth-two MZVs and one depth-one MZV:
\begin{equation}
  \zeta(a)\,\zeta(b)=\zeta(a,b)+\zeta(b,a)+\zeta(a+b).
  \label{found:eq:stuffle-d1}
\end{equation}
The general rule is the same combinatorial device --- the \emph{quasi-shuffle} of Hoffman --- applied to the index words $(s_1,\dots,s_k)$ and $(t_1,\dots,t_\ell)$: interleave the two words in all order-preserving ways, and additionally allow any aligned pair $s_i,t_j$ to \emph{collide} into a single slot $s_i+t_j$. (Setting $a=b=2$ in \eqref{found:eq:stuffle-d1} gives $\zeta(2)^2=2\,\zeta(2,2)+\zeta(4)$, used below.) A depth-two example of the general rule:
\[
  \zeta(a)\,\zeta(b,c)=\zeta(a,b,c)+\zeta(b,a,c)+\zeta(b,c,a)
  +\zeta(a+b,c)+\zeta(b,a+c),
\]
the first three terms inserting $a$ into the three gaps of $(b,c)$, the last two colliding $a$ with $b$ and with $c$.

\paragraph{(b) The shuffle product.}
The integral representation reads each MZV as a Chen iterated integral over the simplex $1>t_1>\dots>t_w>0$ against the two one-forms $\omega_0=\tfrac{dt}{t}$ and $\omega_1=\tfrac{dt}{1-t}$. Encode an admissible index as a binary word: each $s_i$ becomes $\underbrace{0\cdots0}_{s_i-1}1$, so that, e.g., $\zeta(2)\leftrightarrow 01$, $\zeta(3)\leftrightarrow 001$, $\zeta(2,1)\leftrightarrow 011$, $\zeta(3,1)\leftrightarrow 0011$, $\zeta(2,2)\leftrightarrow 0101$. Admissibility ($s_1\ge2$) is the condition that the word \emph{begins with $0$}; convergence at the inner end requires the word to \emph{end with $1$}. Multiplying two iterated integrals over the same simplex produces the integral over the product region, which decomposes into the $\binom{w_1+w_2}{w_1}$ order-preserving interleavings --- the \emph{shuffle product} $\shuffle$ --- of the two letter words, \emph{without} any collision term. Thus
\begin{equation}
  \zeta(a)\,\zeta(b)=\sum_{\text{words }u\,\in\,(\text{word }a)\,\shuffle\,(\text{word }b)}\zeta(u).
  \label{found:eq:shuffle-gen}
\end{equation}

\subsection*{The double-shuffle relation: $\zeta(3,1)$ in closed form}

The canonical demonstration computes $\zeta(2)^2$ by both products and equates the results. Stuffle, from \eqref{found:eq:stuffle-d1} with $a=b=2$:
\begin{equation}
  \zeta(2)^2=2\,\zeta(2,2)+\zeta(4).
  \label{found:eq:z2sq-stuffle}
\end{equation}
Shuffle: $\zeta(2)$ is the word $01$, and we shuffle $01\shuffle01$. Writing the two factors as $0_a1_a$ and $0_b1_b$ and listing the $\binom{4}{2}=6$ interleavings that preserve $0_a\!\to\!1_a$ and $0_b\!\to\!1_b$:
\[
  0_a0_b1_a1_b,\;\;0_a0_b1_b1_a,\;\;0_b0_a1_a1_b,\;\;0_b0_a1_b1_a,\;\;0_a1_a0_b1_b,\;\;0_b1_b0_a1_a.
\]
Forgetting the colour labels, these are the words $0011$ (four times) and $0101$ (twice). Since $0011\leftrightarrow\zeta(3,1)$ and $0101\leftrightarrow\zeta(2,2)$,
\begin{equation}
  \zeta(2)^2=4\,\zeta(3,1)+2\,\zeta(2,2).
  \label{found:eq:z2sq-shuffle}
\end{equation}
Now subtract \eqref{found:eq:z2sq-stuffle} from \eqref{found:eq:z2sq-shuffle}: the $2\,\zeta(2,2)$ terms cancel and we obtain the \emph{double-shuffle relation}
\begin{equation}
  4\,\zeta(3,1)=\zeta(4)
  \qquad\Longrightarrow\qquad
  \zeta(3,1)=\tfrac14\zeta(4)=\frac{\pi^4}{360}.
  \label{found:eq:z31}
\end{equation}
Feeding this back into the sum-theorem relation \eqref{found:eq:sum-wt4} also pins $\zeta(2,2)=\zeta(4)-\zeta(3,1)=\tfrac34\zeta(4)=\pi^4/120$ --- consistent with the all-2's family below.

\begin{remark}[Regularized double shuffle]
The naive double-shuffle relations above are produced only by admissible words. To capture the boundary case, one regularizes the divergent $\zeta(1)$ (the word $1$). Comparing the two regularizations --- shuffle versus stuffle --- as polynomials in a formal variable $T$ standing for the regularized $\zeta(1)$ yields the \emph{extended} or \emph{regularized double-shuffle relations} of Ihara, Kaneko, and Zagier (2006). These are conjectured to generate \emph{all} $\mathbb{Q}$-linear relations among MZVs.
\end{remark}

\subsection*{Other relation families}

The double-shuffle relations sit inside a larger, heavily overlapping web. We list the principal families, one line each.

\paragraph{Duality.} The integral substitution $t\mapsto1-t$ swaps $\omega_0\leftrightarrow\omega_1$ and reverses the order of integration, giving $\zeta(\mathbf{s})=\zeta(\mathbf{s}^\dagger)$, where the \emph{dual} word $\mathbf{s}^\dagger$ is obtained from the binary word by reading it backwards and exchanging $0\leftrightarrow1$. The simplest instance is $011\mapsto011$ read backwards is $110$, complemented is $001$: thus $\zeta(2,1)=\zeta(3)$ \eqref{found:eq:euler-21} is precisely a self/dual coincidence. Likewise $0011\,(\zeta(3,1))$ is \emph{self}-dual---read backwards and complemented it is again $0011$---so duality says nothing new there; the cleanest nontrivial higher instance is the depth-three identity $\zeta(2,1,1)=\zeta(4)$, from $0111\mapsto0001$.

\paragraph{Hoffman's relations.} Relations obtained by comparing the
regularized product with $\zeta(1)$ provide useful index identities. Their
precise insertion and splitting terms matter; a schematic insertion of a
single $1$ in every position is not a valid general formula.

\paragraph{Derivation relations.} An infinite family produced by certain Hopf-algebra derivations acting on the algebra of admissible words (Ihara--Kaneko--Zagier), equivalent over $\mathbb{Q}$ to large portions of the regularized double shuffle.

\paragraph{Ohno's relation.} A single generating identity (Ohno 1999) that simultaneously generalizes the sum theorem \eqref{found:eq:sum-thm} \emph{and} duality, obtained by summing MZVs over all ways of distributing a fixed total increment among the indices of $\mathbf{s}$ and of its dual.

\paragraph{Cyclic sum formula.} Hoffman--Ohno (and Ohno--Wakabayashi): a cyclic sum of MZVs of weight $w$ equals a weighted sum of weight-$w$ values, refining the sum theorem with a cyclic symmetry on the index word.

\subsection*{Closed-form families}

Beyond individual reductions, several infinite families of MZVs evaluate to powers of $\pi$ in closed form. Throughout, $\{s\}_n$ denotes the index $s$ repeated $n$ times.

\paragraph{All 2's.} The generating function $\prod_{n\ge1}(1-x^2/n^2)=\sin(\pi x)/\pi x$ yields, after extracting the relevant power sums,
\begin{equation}
  \zeta(\{2\}_n)=\zeta(\underbrace{2,2,\dots,2}_{n})=\frac{\pi^{2n}}{(2n+1)!}.
  \label{found:eq:all2}
\end{equation}
For $n=1$ this is $\zeta(2)=\pi^2/6$; for $n=2$, $\zeta(2,2)=\pi^4/120$, matching \eqref{found:eq:z31}; for $n=3$, $\zeta(2,2,2)=\pi^6/5040$.

\paragraph{The $\{3,1\}$ family.} A celebrated result of Borwein, Bradley, Broadhurst, and Lison\v{e}k (1998), originally conjectured experimentally:
\begin{equation}
  \zeta(\{3,1\}_n)=\zeta(\underbrace{3,1,3,1,\dots,3,1}_{n\text{ blocks}})=\frac{2\,\pi^{4n}}{(4n+2)!}.
  \label{found:eq:bbbl31}
\end{equation}
For $n=1$ this reproduces $\zeta(3,1)=2\pi^4/6!=\pi^4/360$ of \eqref{found:eq:z31}; for $n=2$, $\zeta(3,1,3,1)=2\pi^8/10!$. The proof is a striking generating-function argument identifying $\sum_n\zeta(\{3,1\}_n)x^{4n}$ with a ratio of $\cosh$ and $\cos$.

\paragraph{The $\{2\}_a\,3\,\{2\}_b$ family.} Zagier (2012) evaluated the two-parameter family $\zeta(\{2\}_a,3,\{2\}_b)$ as an explicit $\mathbb{Q}$-linear combination of products $\pi^{2m}\zeta(2n+1)$, a result whose proof unexpectedly drew on the arithmetic of the modular-form-like generating series; it remains one of the deepest known closed-form families.

\subsection*{Dimensions and the weight grading}

Let $D_w$ be the $\mathbb{Q}$-vector space spanned by all MZVs of weight $w$ (with $D_0=\mathbb{Q}$, $D_1=0$ since $\zeta(1)$ diverges), and $d_w=\dim_\mathbb{Q}D_w$. Zagier (1994) conjectured the remarkably simple recursion
\begin{equation}
  d_w=d_{w-2}+d_{w-3},\qquad d_0=1,\;d_1=0,\;d_2=1,
  \label{found:eq:zagier-dim}
\end{equation}
giving the sequence
\[
  d_0,d_1,d_2,\dots=1,0,1,1,1,2,2,3,4,5,7,9,\dots .
\]
The \emph{upper} bound $d_w\le$ (the recursion value) is a theorem, proved independently by Goncharov and by Terasoma (around 2002), with the conceptual framework of mixed Tate motives completed by Deligne and Goncharov; it amounts to showing that the known relations (e.g.\ associator / motivic) suffice to cut the dimension down. Broadhurst and Kreimer (1997) refined the count by also grading by depth, conjecturing a two-variable generating function that brings in the dimensions of spaces of cusp forms (the first modular obstruction appearing at weight $12$, depth $4$). The matching \emph{lower} bound --- equivalently, that no further relations exist and the conjectured spanning sets are linearly independent --- is wide open; even $\dim D_w\ge2$ for some specific large $w$ would already imply transcendence statements (such as the linear independence of $\pi^6$ and $\zeta(3)^2$ over $\mathbb{Q}$ at $w=6$) that are presently out of reach.

\subsection*{A table of low-weight reductions}

The following table records all admissible MZVs through weight $5$ together with their reductions to the conjectural basis $\{1\}$ (weight $\le2$), $\{\zeta(3)\}$, $\{\zeta(4)\}\sim\{\pi^4\}$, $\{\zeta(5),\zeta(2)\zeta(3)\}$. Every entry has been verified symbolically.

\begin{center}
\renewcommand{\arraystretch}{1.25}
\begin{tabular}{c c l}
\hline
$w$ & spanning count & MZVs of weight $w$ and their reductions\\
\hline
$2$ & $1$ & $\zeta(2)=\dfrac{\pi^2}{6}$\\
$3$ & $1$ & $\zeta(3)$; \quad $\zeta(2,1)=\zeta(3)$\\
$4$ & $1$ & $\zeta(4)=\dfrac{\pi^4}{90}$; \quad $\zeta(3,1)=\dfrac{\pi^4}{360}$, \quad $\zeta(2,2)=\dfrac{\pi^4}{120}$, \quad $\zeta(2,1,1)=\zeta(4)$\\
$5$ & $2$ & $\zeta(5),\ \zeta(2)\zeta(3)$ span; \quad $\zeta(4,1)=2\zeta(5)-\zeta(2)\zeta(3)$,\\
    &     & $\zeta(3,2)=-\tfrac{11}{2}\zeta(5)+3\,\zeta(2)\zeta(3)$, \quad $\zeta(2,3)=\tfrac{9}{2}\zeta(5)-2\,\zeta(2)\zeta(3)$,\\
    &     & $\zeta(3,1,1)=2\zeta(5)-\zeta(2)\zeta(3)$, \quad $\zeta(2,1,1,1)=\zeta(5)$ (duality of $\zeta(5)$)\\
\hline
\end{tabular}
\end{center}

\noindent The table gives spanning sets of sizes $1,1,1,2$ at weights
$2,3,4,5$. The two-element weight-five set is not known to be independent.
Every depth-two value in that row reduces to products of single zeta values.
At weight eight, values such as $\zeta(6,2)$ are candidates for a primitive
direction beyond the polynomial algebra of single zetas.


\section{Harmonic polylogarithms}

The multiple polylogarithms of the previous sections are nested \emph{sums}. For the practitioner who must actually manipulate weight-$w$ transcendentals---most insistently the particle physicist computing Feynman integrals---it is the \emph{iterated-integral} avatar that is decisive, because integration variables, differential equations, and analytic continuation all speak the language of integrals. The harmonic polylogarithms (HPLs) of Remiddi and Vermaseren~\cite{RemiddiVermaseren} are precisely this avatar, specialised to the three integration kernels that arise when a rational integrand has singularities only at $t=0,1,-1$. They have become the de facto standard basis for one-scale problems, and Wolfram~15 exposes them directly as \texttt{HarmonicPolyLog}.

\subsection{Definition: kernels and iterated integrals}

Fix the alphabet $\{0,1,-1\}$ and the three rational \emph{kernels}
\[
  f_{0}(t)=\frac1t,\qquad f_{1}(t)=\frac1{1-t},\qquad f_{-1}(t)=\frac1{1+t}.
\]
A harmonic polylogarithm $H(a_1,\dots,a_w;z)$ is indexed by a word $\vec a=(a_1,\dots,a_w)$ over this alphabet; the length $w$ is the \emph{weight}. The definition is recursive in the leftmost letter. When the leading letter is nonzero,
\begin{equation}\label{found:eq:hpl-rec}
  H(a_1,a_2,\dots,a_w;z)=\int_0^z f_{a_1}(t)\,H(a_2,\dots,a_w;t)\,dt,\qquad a_1\neq 0,
\end{equation}
with the depth-zero seed $H(;z)=1$ (the empty word). The recursion also applies when $a_1=0$, provided the remaining word
is not all zeros. Such an integrand is integrable at the origin. Only the
all-zero words need the logarithmic regularization above; zeros at the end
of a general word can be removed by shuffle identities. For example
$H(0,1;z)=\Li_2(z)$ is a convergent integral with a leading zero.
The leftmost letter always labels the outermost integration.

\subsection{The weight-one and weight-two cases}

At weight one \eqref{found:eq:hpl-rec}--the all-zero regularization give the three building blocks
\begin{equation}\label{found:eq:hpl-w1}
  H(0;z)=\ln z,\qquad H(1;z)=-\ln(1-z),\qquad H(-1;z)=\ln(1+z),
\end{equation}
since $\int_0^z dt/(1-t)=-\ln(1-z)$ and $\int_0^z dt/(1+t)=\ln(1+z)$. At weight two, integrating these against another kernel reproduces the classical dilogarithm and its companions:
\begin{align}
  H(0,1;z)&=\int_0^z\frac{dt}{t}\bigl(-\ln(1-t)\bigr)=\Li_2(z), &
  H(0,-1;z)&=\int_0^z\frac{dt}{t}\ln(1+t)=-\Li_2(-z),\label{found:eq:hpl-w2a}\\
  H(1,0;z)&=\int_0^z\frac{\ln t}{1-t}\,dt=-\Li_2(z)-\ln z\,\ln(1-z), &
  H(1,1;z)&=\tfrac12\ln^2(1-z),\label{found:eq:hpl-w2b}
\end{align}
the last by \eqref{found:eq:hpl-rec} with $H(1;t)=-\ln(1-t)$ and $\int_0^z(-\ln(1-t))/(1-t)\,dt=\tfrac12\ln^2(1-z)$. The pair $H(0,1;z)=\Li_2(z)$ and $H(0,-1;z)=-\Li_2(-z)$ shows the general pattern: the words $H(\underbrace{0,\dots,0}_{s-1},1;z)=\Li_s(z)$ and $H(\underbrace{0,\dots,0}_{s-1},-1;z)=-\Li_s(-z)$ recover the entire ladder of classical polylogarithms, so the HPLs are a strict superset of the functions of Lewin's treatise. In Wolfram~15 these read, e.g., \texttt{HarmonicPolyLog[\{0,1\},x]} $=\Li_2(x)$ and \texttt{HarmonicPolyLog[\{0,-1\},x]} $=-\Li_2(-x)$.

\subsection{The condensed $m$-notation}

Writing out long strings of $0$'s is wasteful, and Remiddi and Vermaseren introduced a condensed index that absorbs them. To each maximal block ``$\underbrace{0,\dots,0}_{|m|-1}\,\sigma$'' consisting of $|m|-1$ zeros followed by a sign letter $\sigma\in\{1,-1\}$ one assigns a single integer
\[
  m=\sigma\cdot(\text{number of zeros}+1),
\]
so the sign of $m$ records $\sigma$ and the magnitude $|m|$ records how many kernels (one $f_{\pm1}$ preceded by $|m|-1$ copies of $f_0$) the block contributes. The resulting tuple is written as a subscript,
\[
  H_{m_1,\dots,m_k}(z):=H(\vec a;z),
\]
and the \emph{depth} of the HPL is $k$, the number of nonzero letters, while the weight is $w=\sum_i|m_i|$. Thus \eqref{found:eq:hpl-w2a} becomes $H_{2}(z)=\Li_2(z)$ and $H_{-2}(z)=-\Li_2(-z)$, the polylogarithm ladder is $H_{s}(z)=\Li_s(z)$, $H_{-s}(z)=-\Li_s(-z)$, and a mixed word such as $H(0,1,-1;z)$ is $H_{2,-1}(z)$. The $m$-notation is the input most reduction tables are organised by; the explicit word notation $H(\vec a;z)$ is the one that makes the integral representation \eqref{found:eq:hpl-rec} transparent. (Some authors reverse the global sign or the letter ordering; we follow the convention compatible with \eqref{found:eq:hpl-w1}--\eqref{found:eq:hpl-w2a} and with \texttt{HarmonicPolyLog}, where indices read outermost-integration-first.)

\subsection{The shuffle algebra}

\begin{theorem}[Shuffle product]\label{found:thm:hpl-shuffle}
For a fixed argument $z$, the product of two harmonic polylogarithms is a $\mathbb{Q}$-linear combination of harmonic polylogarithms of the same argument:
\[
  H(\vec a;z)\,H(\vec b;z)=\sum_{\vec c\in \vec a\shuffle \vec b}H(\vec c;z),
\]
the sum running over all words $\vec c$ in the shuffle $\vec a\shuffle\vec b$, i.e.\ all interleavings of $\vec a$ and $\vec b$ that preserve the internal order of each.
\end{theorem}

This is the iterated-integral analogue of integration by parts: a product of two iterated integrals over $[0,z]$ is the integral over the union of the two simplices, which decomposes into the $\binom{|\vec a|+|\vec b|}{|\vec a|}$ ordered shuffles. Concretely, at weight $1+1$,
\begin{equation}\label{found:eq:hpl-shuffle-ex}
  H(0;z)\,H(1;z)=H(0,1;z)+H(1,0;z),
\end{equation}
which is the statement $\ln z\cdot(-\ln(1-z))=\Li_2(z)+\bigl(-\Li_2(z)-\ln z\,\ln(1-z)\bigr)=-\ln z\,\ln(1-z)$, i.e.\ the integration-by-parts consistency that re-expresses the product of the two weight-one HPLs through the two weight-two ones. (The dilogarithm \emph{reflection} formula is a different relation---the $z\mapsto1-z$ image of $H(0,1;z)$, derived in \eqref{found:eq:dilog-refl} below.) The shuffle algebra is graded by weight, and crucially it lets one strip leading zeros: since $H(0;z)=\ln z$, repeated use of \eqref{found:thm:hpl-shuffle} expresses any $H(0,\vec a;z)$ through $\ln z$ and HPLs whose words begin with a nonzero letter, which is the regularisation alluded to after the all-zero regularization.

\subsection{Values at \texorpdfstring{$z=1$}{z=1}: Euler sums and colored MZVs}

Setting $z=1$ turns admissible HPLs (those whose word does not begin with $1$, so that the boundary integral converges) into constants. Because the kernels live on the fourth-root-of-unity locus $\{0,\pm1\}$, these constants are exactly the level-$2$, or \emph{alternating}, multiple zeta values---the Euler sums. Words over $\{0,1\}$ give the ordinary MZVs,
\begin{equation}\label{found:eq:hpl-mzv}
  H(\underbrace{0,\dots,0}_{s_1-1},1,\underbrace{0,\dots,0}_{s_2-1},1,\dots;1)=\zeta(s_1,s_2,\dots),
\end{equation}
so for instance $H(0,1;1)=\Li_2(1)=\zeta(2)$ and $H(0,0,1;1)=\zeta(3)$. Admitting the letter $-1$ produces the alternating Euler sums: with $\bar s$ marking a sign,
\[
  H(0,-1;1)=-\Li_2(-1)=\tfrac{\pi^2}{12}=\zeta(\bar2),\qquad
  H(0,0,-1;1)=-\Li_3(-1)=\tfrac34\zeta(3),
\]
in the alternating-MZV notation $\zeta(\bar s_1,\dots)=\sum_{n_1>\cdots}\prod(\operatorname{sgn})^{n_i}/n_i^{s_i}$. Thus the $\mathbb{Q}$-algebra of HPL values at $1$ is the algebra of alternating MZVs; its dimensions and the structure of its relations were mapped by Borwein--Bradley--Broadhurst--Lisone\v{k} and, on the motivic side, Deligne and Goncharov. This is the practical reason HPLs are the right basis: the boundary data of a Feynman integral land automatically in a well-understood number-theoretic ring.

\subsection{Argument transformations and the alphabet}
Changing variables transports the singular set of the differential forms.
Negation and inversion, and the map $(1-z)/(1+z)$, preserve the extended
three-letter singular set $\{0,1,-1,\infty\}$, with appropriate endpoint and
branch constants. Reflection $1-z$ preserves the two-letter subalphabet
$\{0,1\}$ but sends the letter $-1$ to $2$. Thus
$H(-1;1-z)=\log(2-z)$ cannot in general be expressed using only the original
three-letter alphabet at $z$. Squaring likewise introduces letters $\pm i$
when applied to a word containing $-1$.

The useful algorithm is to differentiate, partially fraction the transported
kernels, integrate in the resulting alphabet, and fix the endpoint constant.
For the two-letter dilogarithm this gives
\begin{equation}\label{found:eq:dilog-refl}
\Li_2(z)+\Li_2(1-z)=\frac{\pi^2}{6}-\log z\log(1-z),\qquad 0<z<1.
\end{equation}
Analytic continuation requires a common choice of branches.
Duplication is the special closed identity
\[
H(0,1;z^2)=2\bigl(H(0,1;z)-H(0,-1;z)\bigr).
\]
These examples illustrate transport of alphabets, rather than universal
closure under every fractional-linear or quadratic substitution.

\subsection{Relation to Nielsen polylogarithms and to \texorpdfstring{$\Li$}{Li}}

The harmonic polylogarithms contain the Nielsen generalised polylogarithms $S_{n,p}(z)$ of Nielsen (1909), studied analytically by K\"olbig (1986). Indeed $S_{n,p}(z)$ is, exactly the single HPL whose word is $n$ zeros followed by $p$ ones,
\begin{equation}\label{found:eq:nielsen}
  S_{n,p}(z)=H(\underbrace{0,\dots,0}_{n},\underbrace{1,\dots,1}_{p};z),
\end{equation}
of weight $n+p$ and depth $p$; the classical case is $S_{n,1}(z)=\Li_{n+1}(z)$, consistent with \eqref{found:eq:hpl-w2a} for $n=1,p=1$. Thus the chain of inclusions reads: classical $\Li_s$ (depth-$1$ HPLs over $\{0,1\}$) $\subset$ Nielsen $S_{n,p}$ (two-letter words $0^n1^p$) $\subset$ harmonic polylogarithms (arbitrary words over $\{0,1,-1\}$) $\subset$ multiple polylogarithms (the $z=1$ specialisations being MZVs and Euler sums). Each enlargement adds exactly the words the previous class could not name.

\subsection{Finite truncations and Wolfram~15}

Behind every iterated integral stands a nested \emph{sum}, and behind every value at $z=1$ stands the partial sum of that series. These partial sums are the multiple (finite) harmonic numbers
\[
  H^{(s_1,\dots,s_k)}_{N}=\sum_{N\ge n_1>n_2>\cdots>n_k\ge1}\;\prod_{i=1}^{k}\frac{1}{n_i^{s_i}},
\]
generalising the ordinary $H_N=H_N^{(1)}=\sum_{n\le N}1/n$; their $N\to\infty$ limits are the admissible MZVs, and they are the natural objects in which the series expansions of $\varepsilon$-dependent Feynman integrals are stored. Wolfram~15 provides all of these as first-class functions: \texttt{HarmonicPolyLog} for $H(\vec a;z)$ (with the index/sign conventions used above; recall \texttt{HarmonicPolyLog[\{1\},x]} $=-\ln(1-x)$), \texttt{MultiplePolyLog} and \texttt{MultipleZeta} for the limiting transcendentals, and \texttt{MultipleHarmonicNumber} for the finite sums $H^{(\vec s)}_N$. With \texttt{DSolve} and \texttt{AsymptoticDSolveValue} now returning harmonic polylogarithms for Fuchsian systems with rational coefficients, the entire Remiddi--Vermaseren workflow---solve the differential equation, expand at the boundary into Euler sums, continue by the transformations the transformations just discussed---is available without leaving the kernel.

\begin{remark}
The choice of alphabet $\{0,1,-1\}$ is what fixes the singular set at $\{0,1,-1,\infty\}$ and hence ties the boundary values to \emph{alternating} MZVs. Allowing the additional kernels $f_{1/2}(t)=1/(t-\tfrac12)$ or roots of unity of higher order enlarges the class to the two-dimensional harmonic polylogarithms and to Goncharov's multiple polylogarithms with general singular points (the \texttt{GeneralizedPolyLog} of the previous section), at the price of boundary constants drawn from level-$N$ colored MZVs rather than the level-$2$ Euler sums. For one-scale problems the three-letter alphabet is almost always enough, which is the empirical reason HPLs, not their more general cousins, became the working basis of the field.
\end{remark}


\section{Nielsen's generalized polylogarithms}

Long before multiple zeta values acquired their modern algebraic name, Nielsen (1909) introduced a two-parameter family of higher polylogarithms that already pervades Lewin's \emph{Polylogarithms and Associated Functions} (1981, Chapter~7). They are the simplest objects that go beyond the ordinary $\Li_s(z)$ while remaining genuinely ``one-variable'': a single argument $z$, two integer indices $n,p\ge1$, and a closed integral representation. The reader who is comfortable with the dilogarithm reflection and Landen's transformation will find here their direct higher-weight descendants, and -- once we translate into the language of the previous sections -- their place as one distinguished corner of the multiple-polylogarithm world.

\subsection{Definition and elementary reductions}

\begin{definition}
For integers $n,p\ge1$ the \emph{Nielsen generalized polylogarithm} is
\begin{equation}\label{found:eq:nielsen-def}
  S_{n,p}(z)\;=\;\frac{(-1)^{n+p-1}}{(n-1)!\,p!}\int_0^1 \frac{(\ln t)^{n-1}\,\ln^{p}(1-zt)}{t}\,dt ,
\end{equation}
analytically continued from $|z|<1$ to $z\in\mathbb{C}\setminus[1,\infty)$ along the cut of $\ln(1-zt)$. Its \emph{weight} is $n+p$ and its \emph{depth} is $p$.
\end{definition}

The normalizing factorials are chosen exactly so that two boundary specializations collapse onto familiar functions. Setting $p=1$ and integrating $\ln(1-zt)/t$ termwise against $(\ln t)^{n-1}$ recovers the classical Euler integral for the polylogarithm:
\begin{equation}\label{found:eq:Snone}
  S_{n,1}(z)\;=\;\Li_{n+1}(z),\qquad n\ge1 .
\end{equation}
Thus the entire first column of the Nielsen table is nothing but the ordinary polylogarithm shifted by one weight; in particular the lowest member is the dilogarithm,
\begin{equation}\label{found:eq:S11}
  S_{1,1}(z)\;=\;\Li_2(z).
\end{equation}
At the opposite corner $n=1$ the inner factor $\ln t$ disappears and \eqref{found:eq:nielsen-def} reduces to a pure power of $\ln(1-zt)$, giving the elementary
$
  S_{0,p}(z) = \tfrac{(-1)^p}{p!}\,\ln^{p}(1-z)
$
in the degenerate $n=0$ row (we keep $n\ge1$ in the sequel, but this boundary is the seed for the recursions below).

\subsection{Series and the multiple-polylogarithm identification}

Expanding $\ln^{p}(1-zt)$ as a power of the series $-\sum_{m\ge1}(zt)^m/m$ and integrating against $(\ln t)^{n-1}/t$ turns \eqref{found:eq:nielsen-def} into a single power series in $z$ whose coefficients are nested harmonic sums:
\begin{equation}\label{found:eq:nielsen-series}
  S_{n,p}(z)\;=\;\sum_{k\ge1}\frac{z^{k}}{k^{\,n+1}}\;c^{(p)}_{k},
  \qquad
  c^{(p)}_{k}=\sum_{k> m_1>\cdots> m_{p-1}\ge1}\frac{1}{m_1\cdots m_{p-1}} ,
\end{equation}
so $c^{(1)}_k\equiv1$ (reproducing \eqref{found:eq:Snone}) and in general $c^{(p)}_k$ is a depth-$(p-1)$ harmonic sum, an instance of \texttt{MultipleHarmonicNumber}. Comparing \eqref{found:eq:nielsen-series} with the strictly nested sum of the article preamble identifies $S_{n,p}$ as a multiple polylogarithm with a single repeated argument:
\begin{equation}\label{found:eq:nielsen-MPL}
  \boxed{\;S_{n,p}(z)\;=\;\Li_{n+1,\,\underbrace{\scriptstyle 1,\dots,1}_{p-1}}\!\bigl(z,\underbrace{1,\dots,1}_{p-1}\bigr)\;}
\end{equation}
of weight $n+p$ and depth $p$. Its iterated-integral word over the alphabet $\{0,1\}$ (with $0\leftrightarrow dt/t$, $1\leftrightarrow dt/(1-t)$) is
\[
  0^{\,n}\,1^{\,p}\;=\;\underbrace{0\cdots0}_{n}\underbrace{1\cdots1}_{p},
\]
a single block of $n$ zeros followed by a single block of $p$ ones. This is the defining structural fact: among all multiple polylogarithms, the Nielsen family is exactly the ``one initial zero-block, one trailing one-block'' subclass. Everything peculiar about $S_{n,p}$ -- the two-index labelling, the clean ladders, the duality below -- is a shadow of that two-block word.

\begin{remark}[Sign and normalization]
With the positive kernel $dt/(1-t)$ the word $0^n1^p$ equals
$S_{n,p}(z)$ with no extra sign. Replacing each $1$-kernel by
$dt/(t-1)$ gives the factor $(-1)^p$. No factor $(-1)^{p-1}$
belongs to either convention.
\end{remark}

\subsection{Integral and differential relations}

Differentiating \eqref{found:eq:nielsen-series} (or \eqref{found:eq:nielsen-def}) in $z$ lowers the leading index by one and threads the function back into the ladder:
\begin{equation}\label{found:eq:dSnp}
  \frac{d}{dz}\,S_{n,p}(z)\;=\;\frac{1}{z}\,S_{n-1,p}(z)\quad(n\ge1),
  \qquad
  \frac{d}{dz}\,S_{1,p}(z)\;=\;\frac{1}{z}\,S_{0,p}(z)\;=\;\frac{(-1)^{p}}{p\,!}\,\frac{\ln^{p}(1-z)}{z},
\end{equation}
the second relation being the $n=1$ boundary, where the $dt/t$ block is exhausted and the degenerate row $S_{0,p}(z)=\tfrac{(-1)^{p}}{p!}\ln^{p}(1-z)$ is the pure power of $\ln(1-z)$. The two clean ladder moves are the integral forms of these derivatives,
\begin{align}
  \text{zero-block ladder:}\quad
    & S_{n,p}(z)=\int_0^{z}S_{n-1,p}(u)\,\frac{du}{u}\qquad (n\ge1),\label{found:eq:ladder0}\\[2pt]
  \text{one-block ladder:}\quad
    & S_{0,p}(z)=\int_0^{z} S_{0,p-1}(u)\,\frac{du}{1-u}\qquad (p\ge1),\label{found:eq:ladder1}
\end{align}
with seed $S_{0,0}(z)=1$ (the empty word): \eqref{found:eq:ladder0} peels the outermost $0$ from the word $0^{n}1^{p}$, and \eqref{found:eq:ladder1} peels the outermost $1$ from the pure one-block $1^{p}$; in iterated-integral terms both are simply ``integrate against the next letter of $0^{n}1^{p}$.'' The practical content is that the whole Nielsen table is generated from the empty word by first applying the $1$-kernel integral $\int_0^{z}(\cdot)\,du/(1-u)$ a total of $p$ times to build $S_{0,p}$ (so $S_{0,1}(z)=-\ln(1-z)$, $S_{0,2}(z)=\tfrac12\ln^{2}(1-z)$, and so on), then the $0$-kernel integral $\int_0^{z}(\cdot)\,du/u$ a total of $n$ times to prepend the zeros; this is exactly the Fuchsian recursion that \texttt{HarmonicPolyLog} encodes, and it is why \texttt{DSolve} now returns Nielsen polylogs for the relevant rational-coefficient ODEs.

\subsection{Special values}

\subsubsection*{The point $z=1$ and Hoffman's duality}

Setting $z=1$ in \eqref{found:eq:nielsen-MPL} turns the multiple polylogarithm into a multiple zeta value with the same two-block word. The first column reproduces Euler,
\begin{equation}\label{found:eq:Snp-one-col}
  S_{n,1}(1)=\zeta(n+1),\qquad S_{1,2}(1)=\zeta(3),
\end{equation}
and in general
\begin{equation}\label{found:eq:Snp-one}
  S_{n,p}(1)\;=\;\zeta\bigl(n+1,\{1\}^{p-1}\bigr),
\end{equation}
the MZV whose argument string is $n+1$ followed by $p-1$ ones (admissible because the leading entry $n+1\ge2$). Now the iterated-integral word $0^{\,n}1^{\,p}$ of $S_{n,p}$ has a manifest symmetry: reversing it and exchanging $0\leftrightarrow1$ sends $0^{n}1^{p}\mapsto 0^{p}1^{n}$, the word of $S_{p,n}$. This reversal-complement is exactly Hoffman's MZV \emph{duality} (the Kontsevich integral symmetry $\omega\mapsto\overline{\omega}^{\mathrm{rev}}$), so
\begin{equation}\label{found:eq:duality}
  \boxed{\;\zeta\bigl(n+1,\{1\}^{p-1}\bigr)\;=\;\zeta\bigl(p+1,\{1\}^{n-1}\bigr),\qquad\text{equivalently}\qquad S_{n,p}(1)=S_{p,n}(1).\;}
\end{equation}
This is a striking statement: the value of a Nielsen polylogarithm at $1$ is invariant under swapping its two indices, even though $S_{n,p}(z)$ and $S_{p,n}(z)$ are wholly different functions of $z$. For $p=1$ it gives $\zeta(n+1)=\zeta(2,\{1\}^{n-1})$; the first nontrivial case is $S_{1,2}(1)=S_{2,1}(1)$, i.e.\ $\zeta(2,1)=\zeta(3)$, Euler's reflection. We confirmed \eqref{found:eq:Snp-one} and \eqref{found:eq:duality} symbolically against \texttt{MultipleZeta} (e.g.\ $S_{1,2}(1)-\zeta(2,1)$ and $S_{n,p}(1)-S_{p,n}(1)$ vanish to $>70$ digits for all $n,p\le4$).

\subsubsection*{The points $z=-1$ and $z=\tfrac12$ (Kölbig 1986)}

At $z=-1$ the column reduction \eqref{found:eq:Snone} gives $S_{n,1}(-1)=\Li_{n+1}(-1)=-(1-2^{-n})\,\zeta(n+1)$, and the first off-column value is the elegant
\begin{equation}\label{found:eq:S12-minus}
  S_{1,2}(-1)\;=\;\tfrac18\,\zeta(3).
\end{equation}
Higher entries $S_{n,p}(-1)$ reduce (Kölbig 1986) to the alternating-MZV basis at their weight; for example $S_{2,2}(-1)$ and $S_{1,3}(-1)$ already require $\Li_4(\tfrac12)$ alongside $\pi^4$, $\ln^4 2$, $\pi^2\ln^2 2$ and $\zeta(3)\ln 2$, the unmistakable level-$2$ weight-$4$ signature (we recovered both reductions by integer-relation search). At the binary point $z=\tfrac12$ the cleanest closed form is
\begin{equation}\label{found:eq:S12-half}
  \boxed{\;S_{1,2}\!\left(\tfrac12\right)\;=\;\tfrac18\,\zeta(3)\;-\;\tfrac16\,\ln^{3}2.\;}
\end{equation}
We obtained \eqref{found:eq:S12-half} by PSLQ over $\{\zeta(3),\pi^2\ln 2,\ln^3 2\}$ -- the relation $24\,S_{1,2}(\tfrac12)-3\,\zeta(3)+4\ln^3 2=0$ -- and verified it to $40$ digits. (Compare Lewin's dilogarithm value $\Li_2(\tfrac12)=\tfrac{\pi^2}{12}-\tfrac12\ln^2 2$: the same flavour, one weight up.)

\subsection{Functional equations}

The reader who knows Lewin's dilogarithm chapter expects three staples -- reflection $z\leftrightarrow 1-z$, inversion $z\leftrightarrow 1/z$, and the Landen ladder -- to survive at higher weight. They do, as specializations of the harmonic-polylogarithm argument transformations of the previous section, but they grow extra lower-weight ``tail'' terms.

\subsubsection*{Reflection $z\mapsto 1-z$}

At weight $2$ the only reflection is the dilogarithm law itself, here in Nielsen dress ($n=p=1$):
\begin{equation}\label{found:eq:refl-w2}
  S_{1,1}(z)+S_{1,1}(1-z)\;=\;\frac{\pi^2}{6}-\ln z\,\ln(1-z),
\end{equation}
which is exactly Euler's reflection for $\Li_2$. At weight $3$ the genuinely new instance is the $S_{1,2}$ reflection; carrying out the integer-relation search over the natural weight-$3$ basis yields the clean
\begin{equation}\label{found:eq:refl-w3}
  \boxed{\;S_{1,2}(z)\;=\;\zeta(3)\;-\;\Li_3(1-z)\;+\;\ln(1-z)\,\Li_2(1-z)\;+\;\tfrac12\,\ln z\,\ln^{2}(1-z).\;}
\end{equation}
We verified \eqref{found:eq:refl-w3} to $40$ digits at several rational and one negative argument. Two consistency checks are worth noting. Putting $z=1$ collapses the right side to $\zeta(3)-\Li_3(0)+0+0=\zeta(3)$, reproducing \eqref{found:eq:Snp-one-col}; and adding \eqref{found:eq:refl-w3} to its $z\mapsto1-z$ image recovers the symmetric Euler relation $S_{1,2}(z)+S_{1,2}(1-z)=2\zeta(3)-\Li_3(z)-\Li_3(1-z)+\cdots$ with the logarithmic tails arranging themselves into the $\zeta(3)+\Li_3$ shuffle. Identity \eqref{found:eq:refl-w3} is precisely the statement that the iterated integral of $011$ from $0$ to $z$ and from $0$ to $1-z$ differ by an explicit lower-weight boundary correction -- the weight-$3$ avatar of the dilogarithm reflection.

\subsubsection*{Inversion $z\mapsto 1/z$ and Landen}

The inversion does \emph{not} close on $S_{n,p}$ alone at weight $\ge3$: $S_{1,2}(1/z)$ is not a $\mathbb{Q}$-combination of $S_{1,2}(z)$, $\zeta(3)$, $\pi^2\ln z$ and $\ln^3 z$ (an integer-relation search over that ad hoc basis returns only spurious large-coefficient vectors). What is true is the structural statement that $z\mapsto 1/z$, like $z\mapsto1-z$ and the Landen map $z\mapsto z/(z-1)$, is one of the $S_3$ (in fact $\mathfrak{S}$-group) argument transformations under which the \emph{whole} weight-$w$ harmonic-polylogarithm space is closed. Concretely, writing $S_{n,p}$ in the word $0^{n}1^{p}$ and applying the HPL substitution rules of the previous section expresses $S_{n,p}(1/z)$ -- and, separately, $S_{n,p}(z/(z-1))$ -- as an explicit $\mathbb{Q}$-linear combination of all weight-$(n+p)$ harmonic polylogarithms in $z$ together with products of $\ln z$, $\ln(1-z)$ and lower-weight $S$'s and $\zeta$'s. Thus the Nielsen polylogarithms inherit the dilogarithm's full transformation calculus, but only after one leaves the two-index Nielsen subfamily for the ambient HPL algebra; the inversion is the cheapest illustration that $\{S_{n,p}\}$ is \emph{not} closed under the group of fractional-linear argument changes, whereas the HPLs are.

\subsection{Placement inside the multiple-polylogarithm world}

Equation~\eqref{found:eq:nielsen-MPL} identifies $S_{n,p}$ as the ``one trailing block'' corner of the depth-$p$ multiple polylogarithms: the words $0^{n}1^{p}$ form a single two-parameter slice of the full $\{0,1\}^{*}$ alphabet, the slice closest to the ordinary polylogarithm (recovered at $p=1$) and stable under the two operations -- evaluation at $1$ to MZVs, and the duality swap $n\leftrightarrow p$ -- that organize the modern theory. This is exactly why the Nielsen polylogarithms are the natural bridge between Lewin's classical apparatus and the MZV algebra: every Lewin-style identity (reflection, Landen, the special values at $\pm1,\tfrac12$) is the $p\le2$ or $n\le2$ shadow of a multiple-polylogarithm identity, and conversely the duality $S_{n,p}(1)=S_{p,n}(1)$ is the first place a classical reader meets a genuinely MZV-theoretic symmetry with no single-variable analogue. In the Wolfram Language the family is a first-class function: \texttt{PolyLog[n,p,z]} is exactly $S_{n,p}(z)$, with \texttt{PolyLog[n,1,z]} agreeing with \texttt{PolyLog[n+1,z]} per \eqref{found:eq:Snone}, and it sits one specialization step from \texttt{MultiplePolyLog} via \eqref{found:eq:nielsen-MPL} and one from \texttt{MultipleZeta} via \eqref{found:eq:Snp-one}.


\section{The common algebra and the motivic boundary}
The nested series and outer-letter-first integral describe one function.
Choosing the alphabet distinguishes general multiple polylogarithms,
cyclotomic multiple polylogarithms, HPLs and the Nielsen subfamily.
General Goncharov letters can be arbitrary complex numbers; restricting
them to roots of unity gives only the cyclotomic class.

One must distinguish the argument patterns of the two products. Shuffle
at a common endpoint gives $\Li_1(z)^2=2\Li_{1,1}(z,1)$, whereas
stuffle gives $\Li_1(z)^2=2\Li_{1,1}(z,z)+\Li_2(z^2)$.
Both follow directly from their respective representations and therefore
give a nontrivial identity when compared.

Iterated integrals also carry a coproduct organized by subsequences of
marked points. Its reduced components and the associated symbol organize
functional equations by weight. They lose some period constants: vanishing
symbol or reduced coproduct alone is not a proof of a zero function.
Endpoint constants and branches must be fixed separately.
For the dilogarithm the differential
$d\Li_2(z)=-\log(1-z)d\log z$ makes this elementary; differentiating
Euler reflection and evaluating an endpoint proves the identity exactly.

Motivic MZVs~\cite{Brown2012} provide an algebraic setting in which dimension and spanning
theorems can be proved. The numerical period map transports identities
and upper bounds. Its injectivity, and thus the corresponding lower bounds
for numerical periods, is a separate conjectural issue. The same caution
will apply to the Gaussian and Eisenstein depth experiments.
